\documentclass[../main]{subfiles}
\begin{document}
\newpage

\subsubsection{\texttt{oe\_filtered}}

\begin{lstlisting}
theta, m = pysib.oe_filtered(u, y, nb, nf, nz)
\end{lstlisting}

Identifies an OE model using a cost-function shaping strategy to
improve convergence.  The model structure and the criterion minimized
at each stage are the same as for \texttt{oe}:

\[
  y(t) = \frac{B(q^{-1})}{F(q^{-1})}\,u(t - n_k) + e(t), \qquad
  V(\theta) = \frac{1}{N}\sum_t
  \left[y(t) - \frac{B(q^{-1})}{F(q^{-1})}\,u(t-n_k)\right]^2.
\]

Before the final unfiltered OE step, 9 auxiliary OE problems are
solved on filtered versions of the data.  At stage $i$, both
\texttt{u} and \texttt{y} are passed through a first-order low-pass
filter $(1-a_i)/(1-a_i\,q^{-1})$, and the OE criterion is minimized
on the filtered data.  The pole $a_i$ decreases from ${\approx}0.92$
to ${\approx}0.47$, progressively widening the passband from a narrow
low-frequency window toward the full spectrum.  Filtering reshapes the
cost function, reducing the attraction of undesirable local minima
before the estimate is refined on the original data.

\subsubsection*{Arguments}
\begin{center}
\begin{tabular}{lll}
  \toprule
  Name & Type & Description \\
  \midrule
  \texttt{u}  & array\_like & Input signal. \\
  \texttt{y}  & array\_like & Output signal. \\
  \texttt{nb} & int & Number of $B$ parameters. \\
  \texttt{nf} & int & Number of $F$ parameters. \\
  \texttt{nz} & int & Input delay $n_k \geq 0$. \\
  \bottomrule
\end{tabular}
\end{center}

\subsubsection*{Returns}
\begin{center}
\begin{tabular}{lll}
  \toprule
  Name & Type & Value \\
  \midrule
  \texttt{theta} & ndarray & $[b_0,\ldots,b_{n_b-1},\; f_1,\ldots,f_{n_f}]$ \\
  \texttt{m}     & dict    & \texttt{B}, \texttt{F} free; \texttt{A}$=$\texttt{C}$=$\texttt{D}$=[1]$ \\
  \bottomrule
\end{tabular}
\end{center}

\end{document}
